\section*{Hoofdstuk \ref{ch:b1:optical-mineralogy} --- Optische mineralogie}

\begin{solution}{exo:b1:optical-mineralogy:1}
Kwarts: $\qty{30000}{nm}\times0.009 = \qty{270}{nm}$, een \omterm{def:b1:optical-mineralogy:interference-colour}{eerste-orde} wit.
Calciet: $\qty{30000}{nm}\times0.172 = \qty{5160}{nm}$, ver buiten het diagram:
een bleek \omterm{def:b1:optical-mineralogy:interference-colour}{hoge-orde} wit, flauw roze of groen.
\end{solution}

\begin{solution}{exo:b1:optical-mineralogy:2}
$I/I_{\max} = \sin^2 2\theta$: $0$ bij \ang{0}; $\sin^2\ang{45} = 0.5$ bij
\ang{22.5}; $1$ bij \ang{45}; $0.5$ bij \ang{67.5}; $0$ bij \ang{90}.
\end{solution}

\begin{solution}{exo:b1:optical-mineralogy:3}
\omterm{def:b1:rock-forming-minerals:groups}{Granaat} is kubisch, dus isotroop: het licht verlaat het zoals het binnenkwam, langs de
polarisator, en de analysator houdt het tegen; $\Delta = 0$ in elke stand. Een kwartskorrel die
loodrecht op haar \omterm{def:b1:optical-mineralogy:optic-axis}{optische as} is gesneden, heeft ook $n_1 = n_2$ voor licht dat langs die as
reist, dus $\Delta = 0$: ook zij blijft donker. In \omterm{def:b1:field-lab-methods-1:microscope}{vlakgepolariseerd licht} heeft het \omterm{def:b1:rock-forming-minerals:groups}{granaat} zeer
hoog \omterm{def:b1:optical-mineralogy:relief}{reliëf} ($1.83$ tegen $1.54$), de kwarts bijna geen.
\end{solution}

\begin{solution}{exo:b1:optical-mineralogy:4}
$|n - 1.54|$: almandien $0.29$, forsteriet $0.11$, ortoklaas $0.02$, kwarts
$0.01$; \omterm{def:b1:optical-mineralogy:relief}{reliëf} in die volgorde, van hoogste tot laagste. Ortoklaas heeft de lagere
index: negatief \omterm{def:b1:optical-mineralogy:relief}{reliëf}, en haar \omterm{def:b1:optical-mineralogy:relief}{Becke-lijn} beweegt bij het verhogen van de scherpstelling uit de korrel,
het insluitmiddel in.
\end{solution}

\begin{solution}{exo:b1:optical-mineralogy:5}
$\delta = 650/30000 = 0.022$. Een groen \omterm{def:b1:optical-mineralogy:pleochroism}{pleochroïsch} \omterm{def:b1:rock-forming-minerals:mineral}{mineraal} met splijtingen bij
\ang{56} en \ang{124} en een \omterm{def:b1:optical-mineralogy:birefringence}{dubbelbreking} nabij $0.020$: hoornblende.
\end{solution}

\begin{solution}{exo:b1:optical-mineralogy:6}
$\tan\ang{30} = 0.5774$. Kwarts: $(0.013/1.54)\times0.5774 = \qty{0.00487}{rad}
= \ang{0.28}$. Forsteriet: $(0.111/1.54)\times0.5774 = \qty{0.0416}{rad} =
\ang{2.38}$. Verhouding $0.111/0.013 = 8.5$. Exact: $\sin r = 1.54\sin\ang{30}/1.651
= 0.4664$, $r = \ang{27.80}$, een afbuiging van \ang{2.20}: de formule van de eerste \omterm{def:b1:optical-mineralogy:interference-colour}{orde}
overschat met \qty{8}{\percent} voor zo'n groot verschil in index.
\end{solution}

\begin{solution}{exo:b1:optical-mineralogy:7}
$270 + 540 = \qty{810}{nm}$, een \omterm{def:b1:optical-mineralogy:interference-colour}{tweede-orde} groen, wanneer de langzame richtingen van de korrel en het
plaatje evenwijdig zijn (de vertragingen tellen op); $|270 - 540| =
\qty{270}{nm}$, weer \omterm{def:b1:optical-mineralogy:interference-colour}{eerste-orde} wit, wanneer de langzame richting van de korrel evenwijdig is aan de
snelle richting van het plaatje.
\end{solution}

\begin{solution}{exo:b1:optical-mineralogy:8}
Bij volledige dikte is $\Delta = \qty{30000}{nm}\times0.052 = \qty{1560}{nm}$. Vanaf de rand ($\Delta = 0$)
naar binnen passeert de \omterm{def:b1:optical-mineralogy:retardation}{retardatie} 542 en \qty{1074}{nm}, maar niet \qty{1605}{nm}: twee rood-violette randen.
\end{solution}

\begin{solution}{exo:b1:optical-mineralogy:9}
De \omterm{def:b1:optical-mineralogy:extinction}{uitdovingshoek} die op een korrel wordt gemeten, hangt af van hoe de korrel is gesneden: alleen doorsneden in een
bijzondere oriëntatie tonen de grootste hoek, andere tonen minder, tot nul. De grootste waarde, ongeveer \ang{37}, is
de juiste om te vergelijken: zij past bij diopsied (ongeveer \ang{38}), niet bij hoornblende (hoogstens
\ang{19}).
\end{solution}

\begin{solution}{exo:b1:optical-mineralogy:10}
Met de analysator langs $x$ zijn de $x$-componenten van de twee golven
$E_0\cos\theta\cdot\cos\theta$ en $-E_0\sin\theta\,e^{-i\varphi}\cdot(-\sin\theta)$,
dus $A = E_0(\cos^2\theta + \sin^2\theta\,e^{-i\varphi})$ en
\[
  |A|^2 = E_0^2\bigl(\cos^4\theta + \sin^4\theta + 2\sin^2\theta\cos^2\theta\cos\varphi\bigr)
  = E_0^2\bigl(1 - 2\sin^2\theta\cos^2\theta\,(1 - \cos\varphi)\bigr),
\]
met $\cos^4\theta + \sin^4\theta = 1 - 2\sin^2\theta\cos^2\theta$. Omdat
$2\sin^2\theta\cos^2\theta = \frac12\sin^2 2\theta$ en $1 - \cos\varphi =
2\sin^2(\varphi/2)$, is $I_\parallel = I_0(1 - \sin^2 2\theta\sin^2(\pi\Delta/\lambda))$.
Bij elke golflengte geldt $I_\parallel + I_\perp = I_0$: de kleuren tussen evenwijdige
polarisatoren zijn de complementen van die tussen \omterm{def:b1:field-lab-methods-1:microscope}{gekruiste polarisatoren}. Een korrel van
\qty{550}{nm}, magenta tussen \omterm{def:b1:field-lab-methods-1:microscope}{gekruiste polarisatoren}, is tussen evenwijdige polarisatoren bleek geelgroen.
\end{solution}

\begin{solution}{exo:b1:optical-mineralogy:11}
(a) $\Delta/\lambda$ loopt van $5000/700 = 7.14$ tot $5000/400 = 12.5$; het passeert de gehele getallen 8 tot 12:
vijf nullen. (b) De nullen liggen bij $\lambda_k = \Delta/k$; twee opeenvolgende liggen $\Delta/k - \Delta/(k + 1) =
\Delta/(k(k + 1)) \approx \lambda^2/\Delta$ uit elkaar, dat is $550^2/5000 =
\qty{60.5}{nm}$ dicht bij \qty{550}{nm}. (c) Vijf smalle donkere banden verspreid over het spectrum halen een weinig van
elke kleur weg, en het overige wordt doorgelaten met een gemiddeld gewicht van de helft: het oog ziet een bleke, bijna
neutrale tint.
\end{solution}

\begin{solution}{exo:b1:optical-mineralogy:12}
Het gemiddelde van $\sin^2 2\theta$ over een omwenteling is $\frac1{2\pi}\int_0^{2\pi}\sin^2 2\theta\,
\mathrm d\theta = \frac12$, dus de gemiddelde intensiteit is
$\frac12 I_0\sin^2(\pi\Delta/\lambda)$. Zij rekent $\delta = \Delta/30000$
in plaats van $\Delta/25000$ (in nm): al haar waarden zijn $25/30 = 0.83$ van de
juiste, te laag met \qty{17}{\percent}.
\end{solution}

\begin{solution}{pb:b1:optical-mineralogy:1}
\textbf{1.} B: donkere randen betekenen hoog \omterm{def:b1:optical-mineralogy:relief}{reliëf}; de nauwelijks zichtbare grenzen van A betekenen laag \omterm{def:b1:optical-mineralogy:relief}{reliëf}.

\textbf{2.} A: $(1.544 + 1.553)/2 = 1.5485$, $|n - n_m| = 0.0085$. B:
$(1.664 + 1.672 + 1.694)/3 = 1.6767$, $|n - n_m| = 0.1367$.

\textbf{3.} De afbuigingen zijn evenredig met $|n - n_m|$: $0.1367/0.0085 =
16.1$.

\textbf{4.} B heeft een hogere index dan het insluitmiddel: positief \omterm{def:b1:optical-mineralogy:relief}{reliëf}.

\textbf{5.} \omterm{def:b1:optical-mineralogy:pleochroism}{Pleochroïsme}; biotiet.

\textbf{6.} Twee splijtingen bij \ang{56} en \ang{124}: een \omterm{def:b1:rock-forming-minerals:groups}{amfibool}.

\textbf{7.} $\qty{30000}{nm}\times0.009 = \qty{270}{nm}$: \omterm{def:b1:optical-mineralogy:interference-colour}{eerste-orde} wit.

\textbf{8.} Elke korrel is in haar eigen oriëntatie gesneden en toont een waarde van $|n_1 - n_2|$ tussen 0 en
$\delta$; alleen korrels evenwijdig aan de \omterm{def:b1:optical-mineralogy:optic-axis}{optische as} tonen $\delta$ zelf. De hoogste kleur onder vele korrels is
de juiste om te gebruiken.

\textbf{9.} $d = 342/0.009 = \qty{38000}{nm} = \qty{38}{\micro m}$.

\textbf{10.} $(38 - 30)/30 = \qty{26.7}{\percent}$.

\textbf{11.} $\sin^2(\pi\times342/550) = 0.861$; bij \qty{30}{\micro m} is
$\sin^2(\pi\times270/550) = 0.999$.

\textbf{12.} $\qty{38000}{nm}\times0.008 = \qty{304}{nm}$ tot
$\qty{38000}{nm}\times0.013 = \qty{494}{nm}$: wit tot oranjerood, alle onder het einde van de eerste \omterm{def:b1:optical-mineralogy:interference-colour}{orde} bij \qty{542}{nm}.

\textbf{13.} Tussen 1074 en \qty{1605}{nm}: derde \omterm{def:b1:optical-mineralogy:interference-colour}{orde}.

\textbf{14.} $\delta = 1140/38000 = 0.030$.

\textbf{15.} $1140/30000 = 0.038$, wat op een \omterm{def:b1:rock-forming-minerals:groups}{olivijn} zou wijzen. Een magnesiumrijke \omterm{def:b1:rock-forming-minerals:groups}{olivijn} en kwarts reageren tot \omterm{def:b1:rock-forming-minerals:groups}{pyroxeen},
\ce{Mg2SiO4 + SiO2 -> 2 MgSiO3}; een gesteente dat met kwarts is gekristalliseerd, mag beide niet behouden.

\textbf{16.} $0.030$ is de \omterm{def:b1:optical-mineralogy:birefringence}{dubbelbreking} van diopsied; haar hoog \omterm{def:b1:optical-mineralogy:relief}{reliëf}, haar gebrek aan kleur en haar twee splijtingen
nabij \ang{87} stemmen overeen: B is een clinopyroxeen, diopsied.

\textbf{17.} $\qty{30000}{nm}\times0.030 = \qty{900}{nm}$: \omterm{def:b1:optical-mineralogy:interference-colour}{tweede-orde} geel.

\textbf{18.} Met \qty{8}{\micro m}; kwarts toont dan hoogstens \qty{270}{nm}, \omterm{def:b1:optical-mineralogy:interference-colour}{eerste-orde} wit.

\textbf{19.} Polysynthetische \omterm{def:b1:optical-mineralogy:twin}{tweelingvorming}: D is \omterm{def:b1:rock-forming-minerals:groups}{plagioklaas}.

\textbf{20.} De samenstelling van de \omterm{def:b1:rock-forming-minerals:groups}{plagioklaas} (haar anortietgehalte).

\textbf{21.} Een \omterm{def:b1:rock-forming-minerals:groups}{amfibool} met een \omterm{def:b1:optical-mineralogy:extinction}{uitdovingshoek} van \ang{17}: hoornblende (\ang{15}--\ang{19}), geen \omterm{def:b1:rock-forming-minerals:groups}{pyroxeen} (\ang{38} voor diopsied).

\textbf{22.} $\qty{38000}{nm}\times0.020 = \qty{760}{nm}$: een \omterm{def:b1:optical-mineralogy:interference-colour}{tweede-orde} groen.

\textbf{23.} Zij is loodrecht op een \omterm{def:b1:optical-mineralogy:optic-axis}{optische as} gesneden: haar twee indices zijn gelijk voor licht dat langs de as reist,
$\Delta = 0$, en zij blijft donker als een isotrope korrel.

\textbf{24.} A kwarts, B diopsied, C biotiet, D \omterm{def:b1:rock-forming-minerals:groups}{plagioklaas}, E hoornblende: een intermediair stollingsgesteente met wat kwarts,
zoals een kwartsdioriet.

\textbf{25.} De doorsnede is \textbf{\qty[text-series-to-math=true]{38}{\micro m}} dik.
\end{solution}
